% ============================================================
% Capitolo 13 — Comandi conversazionali
% ============================================================

\chapter{Comandi conversazionali}
\label{ch:comandi-conversazionali}

I comandi conversazionali sono istruzioni che dai all'agente nella chat, precedute dal prefisso \texttt{@}. Non sono comandi shell. Sono segnali semantici che l'agente riconosce e interpreta secondo il protocollo ADLC.

\section{La lista completa}
\label{sec:13-lista}

\begin{table}[htbp]
  \centering
  \begin{tabular}{lll}
    \toprule
    \textbf{Comando} & \textbf{Scopo} & \textbf{Modifica file?} \\
    \midrule
    \texttt{@checkpoint} & Fissa stato avanzamento & Propone (no auto) \\
    \texttt{@context-update} & Aggiorna il context & Propone (no auto) \\
    \texttt{@show-constraints} & Mostra vincoli SEC/PERF & No \\
    \texttt{@security-check} & Verifica vs SEC attivi & No \\
    \texttt{@perf-check} & Verifica vs PERF attivi & No \\
    \texttt{@load-phase [N]} & Carica modulo fase N & No \\
    \texttt{@set-phase [N]} & Propone cambio fase & Propone (no auto) \\
    \texttt{@alternatives} & Genera 2--3 opzioni & No (ADR per CRITICAL) \\
    \texttt{@simplify} & Riduce scope & Propone (no auto) \\
    \texttt{@rollback} & Piano di rollback & Propone (no auto) \\
    \texttt{@stop} & Ferma immediatamente & No \\
    \bottomrule
  \end{tabular}
  \caption{Comandi conversazionali ADLC}
  \label{tab:13-comandi}
\end{table}

\section{I comandi più usati in dettaglio}
\label{sec:13-dettaglio}

\subsection*{\texttt{@checkpoint}}

Il più usato. Ogni 3--5 azioni significative o prima di una pausa.

\begin{lstlisting}[language=,caption={Output tipico di @checkpoint}]
Checkpoint — TaskFlow API | 2026-06-27 14:32

Completato:
- POST /tasks: route, service, Zod, RFC 7807, 4 test
- GET /tasks: filtri, paginazione cursor, 6 test

Aggiornamento _CONTEXT.md consigliato:
  Active Task → T-008 GET /tasks/:id
  Active Task Token Est. → 3000/1000/4000
Applico gli aggiornamenti?
\end{lstlisting}

\textbf{Importante:} \texttt{\_CONTEXT.md} e \texttt{PROGRESS.md} non vengono mai modificati automaticamente. L'agente propone il testo; sei tu ad applicarlo.

\subsection*{\texttt{@security-check}}

Revisione del codice/piano contro i vincoli SEC attivi. Ogni claim include il confidence tag.

\begin{lstlisting}[language=,caption={Output di @security-check su POST /tasks}]
SEC-01 Input Validation: [OK]
  Zod schema con titolo obbligatorio, dueDate ISO 8601. [FACT]

SEC-02 Authentication: [OK]
  JWT verificato dal middleware auth. [FACT]

SEC-05 Data Protection: [!] DA VERIFICARE
  Non ho verificato che il log non esponga il body completo.
  [ASSUMPTION — TODO: verifica src/app.ts configurazione Pino]

Azione richiesta: verifica SEC-05 prima del merge.
\end{lstlisting}

\subsection*{\texttt{@alternatives}}

Genera opzioni strutturate con pro, contro e raccomandazione.

\subsection*{\texttt{@simplify}}

Quando il task diventa troppo grande per una sessione: produce uno split in sotto-task con motivazione.

\section{Regole che i comandi non possono violare}
\label{sec:13-regole}

\begin{enumerate}
  \item Nessun comando bypassa gli HALT trigger.
  \item Nessun comando modifica file automaticamente.
  \item \texttt{@security-check} cita solo i vincoli in \texttt{\_CONTEXT.md}.
  \item \texttt{@rollback} non esegue azioni distruttive senza conferma.
  \item I comandi non cambiano la classificazione del rischio.
\end{enumerate}

\section{Comandi personalizzati di progetto}
\label{sec:13-personalizzati}

In \texttt{.adlc/project/instructions.md}:

\begin{lstlisting}[language=,caption={Comandi personalizzati per TaskFlow API}]
## Project Commands

@deploy-check
  Esegui in sequenza:
  1. @security-check sull'endpoint modificato
  2. Verifica che tutti i test passino
  3. Controlla backward-compatibility della migration
  4. Proponi la voce di CHANGELOG.md
\end{lstlisting}

\section*{\LabelSummary}

\begin{itemize}
  \item I comandi \texttt{@} sono segnali semantici --- non comandi shell.
  \item I più usati: \texttt{@checkpoint}, \texttt{@context-update}, \texttt{@security-check}, \texttt{@alternatives}, \texttt{@simplify}, \texttt{@stop}.
  \item Non bypassano il protocollo di sicurezza; non modificano file automaticamente.
  \item Comandi personalizzati in \texttt{.adlc/project/instructions.md} trasformano workflow ricorrenti in shortcut.
\end{itemize}
